\documentclass[12pt,a4paper,oneside]{book}

\usepackage[T1]{fontenc}
\usepackage{lmodern}
\usepackage[margin=2.5cm]{geometry}
\setlength{\headheight}{15pt}
\usepackage{amsmath}
\usepackage{amssymb}
\usepackage{graphicx}
\usepackage{fancyhdr}
\usepackage{hyperref}
\usepackage{fancyvrb}
\usepackage{setspace}
\usepackage{array}
\usepackage{booktabs}
\usepackage{multirow}
\usepackage{float}
\usepackage{caption}
\usepackage{subcaption}


% Python code highlighting
    language=Python,
    escapechar=\,
    basicstyle=\ttfamily\small,
    breaklines=true,
    keywordstyle=\color{blue},
    commentstyle=\color{gray},
    stringstyle=\color{red},
    showstringspaces=false,
    frame=single,
    rulecolor=\color{black!30}
}

% Document metadata
\title{\Large \textbf{ECG Simulator}\\
        \large Comprehensive Technical Documentation\\
        \normalsize GUI, Parameters, and Pathologies}
\author{Enrico Caruso\\Engineering \& Simulation}
\date{\today}

% Header and footer
\pagestyle{fancy}
\fancyhf{}
\fancyhead[L]{\small ECG Simulator Documentation}
\fancyhead[R]{\small Part 2 of 3}
\fancyfoot[C]{\thepage}

% Custom commands
\newcommand{\class}[1]{\texttt{#1}}
\newcommand{\func}[1]{\texttt{#1}()}
\newcommand{\param}[1]{\textit{#1}}

\onehalfspacing

\begin{document}

\frontmatter

\tableofcontents
\listoffigures
\listoftables

\mainmatter

% ================================================================
\chapter{GUI Description: Complete Reference}
% ================================================================

File: \texttt{cardiac\_sim/gui/main\_window.py}

\section{Main Window Architecture}

The application window has three main regions:

\begin{figure}[H]
    \centering
    
    \begin{center}
    \texttt{[Diagram - text description in document]}
    \end{center}
    
    \caption{Main window layout. Toolbar at top; left sidebar with parameter/pathology panels; central area for ECG display.}
    \label{fig:main_window}
\end{figure}

\subsection{Toolbar}

Located at the top, the toolbar provides:

\begin{enumerate}
    \item \textbf{Run/Pause button}: Toggle simulation on/off
    \begin{itemize}
        \item Enabled: simulation running, ECG updating continuously
        \item Disabled: simulation paused, ECG frozen
        \item Default: running
    \end{itemize}
    
    \item \textbf{Reset button}: Clear ECG history and restart simulation from $t=0$
    \begin{itemize}
        \item Clears all buffers
        \item Reinitializes engine
        \item Restarts timer
    \end{itemize}
    
    \item \textbf{Engine selector combo box}: Choose between \phaseI and \phaseIII
    \begin{itemize}
        \item Option 1: ``Phase 1/2 – Conduction Graph'' (default)
        \item Option 2: ``Phase 3 – FHN Cable''
        \item Selecting triggers engine swap with no parameter loss
        \item Implementation: \texttt{\_on\_engine\_changed} in main\_window.py
    \end{itemize}
    
    \item \textbf{Heart rate display}: Reads parameters from engine and shows current HR
    \begin{itemize}
        \item Format: ``HR: XX bpm''
        \item Updated every GUI frame (typically 30--60 Hz)
    \end{itemize}
\end{enumerate}

\section{Left Sidebar: Parameter Panel}

File: \texttt{cardiac\_sim/gui/parameter\_panel.py}

\subsection{Overview}

The left sidebar contains a hierarchical parameter tree. Each parameter is a \textbf{named value} with:

\begin{itemize}
    \item \textbf{Name}: descriptive label
    \item \textbf{Type}: numeric (int, float), boolean, or choice
    \item \textbf{Default value}: physiological normal
    \item \textbf{Minimum/maximum}: physiologically plausible range
    \item \textbf{Units}: seconds, milliseconds, dimensionless, etc.
    \item \textbf{Current value}: user-adjustable
\end{itemize}

The tree is organized into categories (sections) for clarity:

\begin{Verbatim}[formatcom=\color{black}]
Parameters:
+- SA Node
|  +- Cycle Length (ms)
|  +- Ectopic Trigger Rate (beats/min)
|  └- [other SA parameters]
+- AV Node
|  +- AV Delay (ms)
|  └- [others]
+- Refractory Periods
|  +- Atrial (ms)
|  +- Ventricular (ms)
|  └- [others]
+- Phase 3 Model (if Phase 3 selected)
|  +- Cable Diffusion Atrial (units)
|  +- Cable Diffusion Ventricular (units)
|  └- Time Scale Tau (ms)
└- [other categories]
\end{Verbatim}

\subsection{Complete Parameter List with Ranges and Meanings}

\begin{center}
\begin{longtable}{|l|c|c|c|p{3cm}|}
\hline
\textbf{Parameter Name} & \textbf{Min} & \textbf{Default} & \textbf{Max} & \textbf{Meaning} \\
\hline
\endhead

\hline
\endfoot

\multicolumn{5}{|c|}{\textbf{SA Node Section}} \\
\hline
Cycle Length (ms) & 300 & 857 & 2000 & Time between SA fires; inverse of heart rate \\
\hline
Ectopic Trigger Rate (bpm) & 0 & 0 & 200 & Ectopic atrial pacemaker rate if enabled \\
\hline
\multicolumn{5}{|c|}{\textbf{Atrial Section}} \\
\hline
RA Conduction Delay (ms) & 0 & 50 & 200 & Extra delay in right atrium \\
\hline
LA Conduction Delay (ms) & 0 & 50 & 200 & Extra delay in left atrium \\
\hline
RA–LA Conduction Delay (ms) & 0 & 100 & 300 & RA to LA propagation time \\
\hline
\multicolumn{5}{|c|}{\textbf{AV Node Section}} \\
\hline
AV Delay (ms) & 50 & 100 & 300 & Programmed AV node delay \\
\hline
AV Conductance & 0 & 1.0 & 1.0 & 0 = complete block; 1 = normal conduction \\
\hline
AV Refractory Period (ms) & 200 & 300 & 400 & Minimum interval before re-excitation \\
\hline
\multicolumn{5}{|c|}{\textbf{Ventricular Section}} \\
\hline
RV Activation Delay (ms) & 0 & 20 & 100 & Extra delay in right ventricle \\
\hline
LV Activation Delay (ms) & 0 & 30 & 100 & Extra delay in left ventricle \\
\hline
SEPT Activation Delay (ms) & 0 & 10 & 50 & Septal activation delay \\
\hline
Septal Conductance & 0 & 1.0 & 1.0 & 0 = block septal conduction \\
\hline
RBB Conductance & 0 & 1.0 & 1.0 & 0 = complete RBBB; 1 = normal \\
\hline
LBB Conductance & 0 & 1.0 & 1.0 & 0 = complete LBBB; 1 = normal \\
\hline
Ventricular Refractory Period (ms) & 200 & 250 & 350 & Ventricular refractoriness \\
\hline
\multicolumn{5}{|c|}{\textbf{Phase 3 Model (if enabled)}} \\
\hline
Cable Time Scale $\tau$ (ms) & 5 & 15 & 30 & FHN time constant; affects overall speed \\
\hline
Atrial Diffusion Coefficient & 0.5 & 2.88 & 10.0 & Atrial cell-to-cell coupling \\
\hline
Ventricular Diffusion LBB & 50 & 72 & 150 & His–LBB–LV diffusion (fast pathway) \\
\hline
Ventricular Diffusion LV & 1 & 3.0 & 20 & LV myocardial diffusion (slower) \\
\hline
\multicolumn{5}{|c|}{\textbf{ECG Display}} \\
\hline
Display Window (s) & 2 & 10 & 30 & Horizontal axis span \\
\hline
Amplitude Scale (mV/div) & 0.1 & 0.5 & 2.0 & Vertical scale; standard ECG = 0.5 mV/div \\
\hline

\end{longtable}
\end{center}

\subsection{Parameter Data Model}

File: \texttt{cardiac\_sim/core/parameter\_model.py}

\begin{Verbatim}[formatcom=\color{black}]
@dataclass
class SimulationParameters:
    """Central holder of all tunable parameters."""
    
    # SA node
    sa_node: SANodeParams
    
    # Conduction delays
    atrial: AtrialParams
    av_node: AVNodeParams
    ventricular: VentricularParams
    
    # Refractory periods (Phase 1/2)
    refractoriness: RefractoryParams
    
    # Phase 3 parameters
    cable_model: CableModelParams
    
    # ECG display
    ecg_display: ECGDisplayParams
\end{Verbatim}

Each parameter group is a dataclass with validation:

\begin{Verbatim}[formatcom=\color{black}]
@dataclass
class AVNodeParams:
    av_delay_ms: float = 100.0    # Normalized to 100 ms
    conductance: float = 1.0       # 0–1 scale: 0=block, 1=normal
    refractory_ms: float = 300.0
    
    def __post_init__(self):
        """Validate ranges."""
        assert 50 <= self.av_delay_ms <= 300, "AV delay out of range"
        assert 0 <= self.conductance <= 1, "Conductance must be 0–1"
\end{Verbatim}

\section{Left Sidebar: Pathology Panel}

File: \texttt{cardiac\_sim/gui/pathology\_panel.py}

\subsection{Overview}

The Pathology Panel is a \textbf{list of checkboxes}, each representing a pathological condition. When checked, the condition is applied to the current engine.

\subsection{Available Pathologies}

\begin{center}
\begin{longtable}{|l|p{5cm}|}
\hline
\textbf{Pathology Name} & \textbf{Description and Implementation} \\
\hline
\endhead

\hline
\endfoot

\multicolumn{2}{|c|}{\textbf{AV Blocks}} \\
\hline
First-Degree AV Block & Increase AV delay by 50 ms (total ≈ 150 ms). PR interval prolonged; QRS unchanged. \\
\hline
Second-Degree AV Block (Mobitz I) & Progressively increase AV delay; eventually a beat is dropped. Not implemented in current version. \\
\hline
Second-Degree AV Block (Mobitz II) & Randomly drop 1 in $N$ beats from AV conduction. \\
\hline
Third-Degree AV Block (Complete) & Set AV conductance = 0. Ventricular escape rhythm emerges (junctional or idioventricular). \\
\hline
\multicolumn{2}{|c|}{\textbf{Bundle Branch Blocks}} \\
\hline
Right Bundle Branch Block (RBBB) & Set RBB conductance = 0. RV is activated late via collateral LV conduction; QRS widened, characteristic terminal R' in V1. \\
\hline
Left Bundle Branch Block (LBBB) & Set LBB conductance = 0. LV is activated late via collateral RV conduction; QRS widened, broad notched deflection in lateral leads. \\
\hline
\multicolumn{2}{|c|}{\textbf{Atrial Pathologies}} \\
\hline
Atrial Fibrillation & Replace orderly BFS with chaotic re-entrant activation. Atrial rate 300–600 bpm; erratic baseline; absent P waves. \\
\hline
Atrial Flutter & Regular re-entrant circuit; atrial rate 250–350 bpm; regular ``sawtooth'' pattern visible in II, III, aVF. \\
\hline
Sinus Bradycardia & Increase SA cycle length to 1200 ms (50 bpm). Normal QRS and T morphology; just slower rate. \\
\hline
Sinus Tachycardia & Decrease SA cycle length to 400 ms (150 bpm). Normal morphology; faster rate. \\
\hline
\multicolumn{2}{|c|}{\textbf{Ventricular Pathologies}} \\
\hline
Premature Ventricular Contraction (PVC) & Trigger ectopic ventricular activation at preset time (usually 600 ms post-SA). Produces wide QRS; T-wave opposite polarity. \\
\hline
Ventricular Tachycardia & Repetitive PVCs at rate 120–250 bpm. Rapid wide complexes; may degenerate to VF. \\
\hline
Ventricular Fibrillation & Chaotic ventricular activation; no organized QRS; baseline appears ``fine'' or ``coarse'' fibrillary waves. \\
\hline
\multicolumn{2}{|c|}{\textbf{Accessory Pathways}} \\
\hline
Wolff-Parkinson-White (WPW) & Add accessory pathway bypassing AV node; early ventricular pre-excitation creates short PR and delta wave. \\
\hline
Atrioventricular Reentrant Tachycardia (AVNRT) & Dual AV node pathways; setup conditions for re-entrant circuit; produces regular narrow-complex tachycardia. \\
\hline

\end{longtable}
\end{center}

\subsection{Pathology Plugin System}

File: \texttt{cardiac\_sim/pathology/\_\_init\_\_.py} and \texttt{cardiac\_sim/plugins/}

All pathologies are implemented as \textbf{plugins}---modular classes inheriting from an abstract base:

\begin{Verbatim}[formatcom=\color{black}]
class AbstractPathologyPlugin(ABC):
    """Base class for all pathologies."""
    
    @property
    @abstractmethod
    def name(self) -> str:
        """Pathology name for UI."""
        pass
    
    @abstractmethod
    def apply(self, engine: SimulationEngine, params: SimulationParameters) -> None:
        """Apply pathological modifications to engine and/or parameters."""
        pass
    
    @abstractmethod
    def unapply(self, engine: SimulationEngine, params: SimulationParameters) -> None:
        """Restore normal state."""
        pass
\end{Verbatim}

Example implementation (RBBB):

\begin{Verbatim}[formatcom=\color{black}]
# File: cardiac_sim/plugins/bundle_blocks.py

class RBBBPathology(AbstractPathologyPlugin):
    @property
    def name(self) -> str:
        return "Right Bundle Branch Block"
    
    def apply(self, engine: SimulationEngine, params: SimulationParameters) -> None:
        # Store original for later restoration
        self._original_rbb_conductance = params.ventricular.rbb_conductance
        
        # Apply block
        params.ventricular.rbb_conductance = 0.0
        
        # Some engines may need explicit reinitialization
        engine.apply_pathology(self)
    
    def unapply(self, engine: SimulationEngine, params: SimulationParameters) -> None:
        # Restore
        params.ventricular.rbb_conductance = self._original_rbb_conductance
        engine.apply_pathology(None)  # Clear pathology
\end{Verbatim}

\subsection{Pathology GUI Interaction}

When the user checks a pathology checkbox:

\begin{enumerate}
    \item GUI calls \texttt{apply} on the plugin
    \item Plugin modifies parameters and/or engine state
    \item Engine is reinitialized with new parameters (seamlessly)
    \item ECG display updates to show new morphology
\end{enumerate}

When unchecked:

\begin{enumerate}
    \item GUI calls \texttt{unapply}
    \item Original parameters restored
    \item Engine reinitialized
    \item ECG returns to normal
\end{enumerate}

\section{Central Display: ECG Plotter}

File: \texttt{cardiac\_sim/gui/ecg\_display.py}

\subsection{Overview}

The ECG display is a custom PyQt6 widget using \textbf{pyqtgraph} for high-performance waveform rendering.

\begin{Verbatim}[formatcom=\color{black}]
class ECGDisplayWidget(pg.GraphicsLayoutWidget):
    """Real-time 12-lead ECG display."""
    
    def __init__(self, parent=None):
        super().__init__(parent)
        
        # Create 12 plot items (one per lead)
        self._plots = {}
        for lead_name in LEAD_NAMES:  # ['I', 'II', ..., 'V6']
            plot = self.addPlot(title=lead_name, axisItems={'bottom': TimeAxis()})
            curve = plot.plot(pen='blue')
            self._plots[lead_name] = {'plot': plot, 'curve': curve, 'data': []}
    
    def update_ecg(self, ecg_sample: ECGSample) -> None:
        """Append new sample and refresh display."""
        for i, lead_name in enumerate(LEAD_NAMES):
            self._plots[lead_name]['data'].append(ecg_sample.leads[i])
            
            # Keep only last N samples (for time window)
            if len(self._plots[lead_name]['data']) > self.max_samples:
                self._plots[lead_name]['data'].pop(0)
            
            # Update plot
            self._plots[lead_name]['curve'].setData(
                self._plots[lead_name]['data'],
                pen=pg.mkPen('blue', width=2)
            )
\end{Verbatim}

\subsection{Rendering Parameters}

\begin{itemize}
    \item \textbf{Display window}: Default 10 seconds; configurable in toolbar
    \item \textbf{Amplitude scale}: Default 0.5 mV/div (clinical standard); configurable
    \item \textbf{Refresh rate}: 30–60 Hz (user-transparent)
    \item \textbf{Grid}: Major gridlines every 200 ms (horizontal) and 0.5 mV (vertical)
    \item \textbf{Color}: Blue curves, black background
\end{itemize}

\subsection{Performance Considerations}

To maintain real-time performance:

\begin{itemize}
    \item Data kept in circular buffers (avoid continuous reallocation)
    \item Rendering occurs on separate Qt thread (GUI remains responsive)
    \item Downsampling applied if necessary (though typically unnecessary)
\end{itemize}

\section{Simulation Worker Thread}

File: \texttt{cardiac\_sim/core/simulation\_worker.py}

\subsection{Purpose}

The simulation engine runs on a separate thread to prevent GUI freezing. The \texttt{SimulationWorker} class encapsulates this threading:

\begin{Verbatim}[formatcom=\color{black}]
class SimulationWorker(QThread):
    """Worker thread for simulation engine."""
    
    sample_generated = pyqtSignal(ECGSample)  # Emitted when ECG sample ready
    error = pyqtSignal(str)
    
    def __init__(self):
        super().__init__()
        self._engine: SimulationEngine = ConductionGraphEngine()
        self._is_running: bool = False
        self._params: SimulationParameters = SimulationParameters()
    
    def run(self) -> None:
        """Main simulation loop (runs on worker thread)."""
        dt = 0.002  # 2 ms per frame (500 Hz simulation rate)
        
        while self._is_running:
            sample = self._engine.step(dt)
            self.sample_generated.emit(sample)
            # Throttle to ~500 Hz
            time.sleep(dt / 2)  # Rough throttling
    
    def set_engine(self, new_engine: SimulationEngine, params: SimulationParameters) -> None:
        """Swap engines (must be called when NOT running)."""
        if self._is_running:
            raise RuntimeError("Cannot swap engines while running")
        
        self._engine = new_engine
        self._params = params
        self._engine.initialize(params)
\end{Verbatim}

\section{Data Flow Diagram}

\begin{figure}[H]
    \centering
    
    \begin{center}
    \texttt{[Diagram - text description in document]}
    \end{center}
    
    \caption{Data flow from UI components through engines to ECG display.}
    \label{fig:dataflow}
\end{figure}

% ================================================================
\chapter{Detailed Implementation: Cable Engine and Cable Model}
% ================================================================

\section{ConductionCableEngine: Complete Details}

File: \texttt{cardiac\_sim/simulation/cable\_engine.py} (≈ 450 lines)

\subsection{Overview}

The \texttt{ConductionCableEngine} wraps a \texttt{ConductionCable1D} numerical model and manages beat state, ECG computation, and multi-beat rhythm.

\subsection{Key Attributes}

\begin{Verbatim}[formatcom=\color{black}]
class ConductionCableEngine(SimulationEngine):
    def __init__(self):
        self._cable: ConductionCable1D = None
        self._params: SimulationParameters = None
        self._nodes: dict[str, ECGNode] = {}
        self._forward_model: LeadFieldModel = None
        
        # Beat management
        self._time: float = 0.0
        self._next_sa_fire: float = 0.0
        self._active_beats: list[BeatRecord] = []  # Up to 3 concurrent beats
        self._beat_in_progress: bool = False
        self._current_beat_nodes: dict[str, float] = {}
        
        # Performance tracking
        self._step_count: int = 0
\end{Verbatim}

\subsection{Initialization: \texttt{initialize}}

\begin{Verbatim}[formatcom=\color{black}]
def initialize(self, params: SimulationParameters) -> None:
    """Create cable and build node dictionary."""
    self._params = params
    self._cable = ConductionCable1D(params)
    self._nodes = _build_cable_nodes(params)  # Adjusted Gaussians
    self._forward_model = LeadFieldModel()
    self._time = 0.0
    self._next_sa_fire = 0.0
    self._active_beats.clear()
\end{Verbatim}

\subsection{Main Simulation Loop: \texttt{step}}

This is the \textbf{critical function} that drives each simulation step:

\begin{Verbatim}[formatcom=\color{black}]
def step(self, dt_s: float) -> ECGSample:
    """Advance by dt_s, return ECG sample.
    
    KEY PRINCIPLE: Always advance cable (never skip based on beat state).
    This ensures internal cable clock is synchronized with wall clock.
    """
    self._time += dt_s
    
    # --- SA pacemaker ------------------------------------
    if self._time >= self._next_sa_fire:
        self._fire_sa()
    
    # --- Always integrate cable (CRITICAL) ----------------
    newly_activated = self._cable.advance(dt_s)
    
    # --- Collect newly activated nodes -------------------
    if self._beat_in_progress:
        for node_name, t_act in newly_activated.items():
            if node_name not in self._current_beat_nodes:
                self._current_beat_nodes[node_name] = t_act
                
                # KEY FIX: update active beat immediately
                # so ECG renders QRS at activation time
                if self._active_beats:
                    self._active_beats[-1].activation_times[node_name] = t_act
    
    # --- Check beat completion --------------------------
    self._check_beat_completion()
    
    # --- Compute ECG ------------------------------------
    ecg = self._compute_ecg(self._time)
    
    return ECGSample(timestamp=self._time, leads=ecg)
\end{Verbatim}

\subsection{SA Pacemaker: \texttt{\_fire\_sa}}

\begin{Verbatim}[formatcom=\color{black}]
def _fire_sa(self) -> None:
    """Initiate a new beat (SA node fires)."""
    cycle_len = self._params.sa_node.cycle_length_ms / 1000.0
    self._next_sa_fire += cycle_len
    
    # Start new beat
    beat_record = BeatRecord(
        beat_number=len(self._active_beats),
        sa_fire_time=self._time,
        activation_times={},  # Will populate as nodes activate
        qrs_start=None,
        qrs_end=None
    )
    self._active_beats.append(beat_record)
    self._beat_in_progress = True
    self._current_beat_nodes.clear()
    
    # Trigger cable to begin new beat
    self._cable.new_beat(self._time, self._params)
\end{Verbatim}

\subsection{Beat Completion: \texttt{\_check\_beat\_completion}}

A beat is considered complete when:
\begin{enumerate}
    \item All ventricular nodes have been activated, OR
    \item Time exceeds $0.9 \times$ cycle length (adaptive timeout)
\end{enumerate}

\begin{Verbatim}[formatcom=\color{black}]
def _check_beat_completion(self) -> None:
    """Commit beat when all nodes activated or timeout reached."""
    if not self._beat_in_progress or not self._active_beats:
        return
    
    beat = self._active_beats[-1]
    cycle_len = self._params.sa_node.cycle_length_ms / 1000.0
    
    # All ventricular nodes activated?
    ven_nodes = set(VENTRICULAR_NODES)  # pre-defined
    activated = set(beat.activation_times.keys())
    ven_complete = ven_nodes.issubset(activated)
    
    # Timeout check
    time_limit = beat.sa_fire_time + 0.9 * cycle_len
    timeout_reached = self._time >= time_limit
    
    if ven_complete or timeout_reached:
        beat.is_committed = True
        self._beat_in_progress = False
        
        # Store QRS timing for diagnostics
        if beat.activation_times:
            times_v = [
                t for n, t in beat.activation_times.items()
                if n in VENTRICULAR_NODES
            ]
            if times_v:
                beat.qrs_start = min(times_v)
                beat.qrs_end = max(times_v)
\end{Verbatim}

\subsection{ECG Computation: \texttt{\_compute\_ecg}}

\begin{Verbatim}[formatcom=\color{black}]
def _compute_ecg(self, t_s: float) -> np.ndarray:
    """Sum Gaussian dipoles from all active beats.
    
    Returns:
        np.ndarray of shape (12,) with voltages in mV
    """
    dipole_3d = np.zeros(3)  # X, Y, Z components
    
    # Retention window: only beats within last 0.85 s contribute
    retention_window = 0.85
    t_cutoff = t_s - retention_window
    
    for beat in self._active_beats:
        if beat.sa_fire_time < t_cutoff:
            # Remove old beats to prevent memory growth
            self._active_beats.remove(beat)
            continue
        
        # Sum contributions from all nodes
        for node_name, t_act in beat.activation_times.items():
            if node_name not in self._nodes:
                continue
            
            node = self._nodes[node_name]
            
            # Depolarization Gaussian
            tau_depol = node.depol_duration  # typically 12.7 ms
            g_depol = np.exp(-(t_s - t_act)**2 / (2 * tau_depol**2))
            if t_s < t_act:
                g_depol = 0
            
            # Repolarization Gaussian
            t_repol = t_act + node.repol_delay  # peak at ~280 ms post-act
            tau_repol = node.repol_duration
            g_repol = np.exp(-(t_s - t_repol)**2 / (2 * tau_repol**2))
            
            # Dipole contribution
            dipole_node = node.dipole_vector * (
                g_depol - node.repol_amplitude * g_repol
            )
            dipole_3d += dipole_node
    
    # Project onto 12 lead vectors
    ecg_12 = self._forward_model.project(dipole_3d)
    return ecg_12
\end{Verbatim}

\section{ConductionCable1D: Numerical Details}

File: \texttt{cardiac\_sim/core/tissue/cable\_1d.py} (≈ 500 lines)

\subsection{Cell Configuration}

\begin{Verbatim}[formatcom=\color{black}]
# Fixed anatomical structure
_NA = 8   # Atrial cells: SA(0) + RA(1-3) + LA(4-7)
_NV = 10  # Ventricular cells: His(0) + LBB(1-3) + LV(4-9)

# Cell segmentation (maps to nodes)
_ATR_SEGS = {
    'SA': (0, 1),           # SA node (index 0)
    'RA_EARLY': (1, 2),     # RA early
    'RA_MID': (2, 3),
    'RA_LATE': (3, 4),
    'LA_EARLY': (4, 5),     # LA early
    'LA_MID': (5, 6),
    'LA_LATE': (6, 8),      # LA late (indices 6-7)
}

_VEN_SEGS = {
    'HIS': (0, 1),
    'LBB_PROX': (1, 2),
    'LBB_MID': (2, 3),
    'LBB_DIST': (3, 4),
    'LV_SEPT': (4, 5),      # LV septum (early)
    'LV_ANT': (5, 6),
    'LV_LAT': (6, 7),       # LV lateral (late, ~645 ms)
    'LV_INF': (7, 8),
    'LV_POST': (8, 9),      # LV posterior
    'RV': (9, 10),          # Actually computed separately; not in sequential cable
}
\end{Verbatim}

\subsection{Diffusion Coefficients (Phase 3-D Calibrated)}

\begin{Verbatim}[formatcom=\color{black}]
# Atrial cable
_D_ATR = np.array([
    0.0,      # SA node (point source, no diffusion)
    2.88,     # RA
    2.88,
    2.88,
    2.88,     # LA
    2.88,
    2.88,
    2.88
])

# Ventricular cable
_D_VEN = np.array([
    0.0,      # His bundle (point source)
    72.0,     # LBB proximal (fast conduction, D_LV=3 earlier made QRS too wide)
    72.0,     # LBB mid
    72.0,     # LBB distal
    3.0,      # LV begins
    3.0,      # LV anterior
    3.0,      # LV lateral
    3.0,      # LV inferior
    3.0,      # LV posterior
    3.0       # RV auxiliary
])

# CFL check at initialization:
D_max_atr = 2.88
D_max_ven = 72.0
dt_int = 0.0001    # 0.1 ms micro-step
tau = 0.015        # 15 ms physical time scale

CFL_atr = D_max_atr * dt_int / (tau * 1**2) = 2.88 * 0.0001 / 0.015 = 0.0192  ✓
CFL_ven = D_max_ven * dt_int / (tau * 1**2) = 72.0 * 0.0001 / 0.015 = 0.480   ✓ (marginal)
\end{Verbatim}

\subsection{FHN State Variables}

\begin{Verbatim}[formatcom=\color{black}]
class ConductionCable1D:
    def __init__(self, params: SimulationParameters):
        # Atrial state: (cells, [voltage, recovery])
        self._sa = np.zeros((_NA, 2))  # Shape: (8, 2)
        self._sv = np.zeros((_NV, 2))  # Shape: (10, 2)
        
        # Activation tracking
        self._t_act_a = np.full(_NA, math.inf)  # activation time per cell
        self._t_act_v = np.full(_NV, math.inf)
        
        # Internal clock
        self._t = 0.0
        
        # Stimulus timing
        self._stim_end_a = 0.0
        self._av_pending_at = math.inf  # when AV should fire
        
        # Pathology state
        self._av_conductance = params.av_node.conductance  # 0–1
        self._sep_conductance = params.ventricular.septal_conductance
        self._rbb_conductance = params.ventricular.rbb_conductance
        self._lbb_conductance = params.ventricular.lbb_conductance
\end{Verbatim}

\subsection{New Beat: \texttt{new\_beat}}

\begin{Verbatim}[formatcom=\color{black}]
def new_beat(self, sa_fire_time: float, params: SimulationParameters) -> None:
    """Reset cable state for new beat.
    
    Initial conditions: all cells at rest, with stimulus applied
    to SA node.
    """
    self._t = sa_fire_time
    
    # Initialization
    V_REST, W_REST = FitzHughNagumoCell.get_resting_state()
    self._sa[:, 0] = V_REST   # voltage
    self._sa[:, 1] = W_REST   # recovery
    self._sv[:, 0] = V_REST
    self._sv[:, 1] = W_REST
    
    # Mark activation times
    self._t_act_a[:] = math.inf
    self._t_act_v[:] = math.inf
    self._t_act_a[0] = sa_fire_time  # SA node fires at designated time
    
    # Stimulus window
    self._stim_end_a = sa_fire_time + _STIM_DUR
    self._av_pending_at = math.inf
    
    # Refractory enforcement
    self._last_activation_v = np.full(_NV, -math.inf)
\end{Verbatim}

\subsection{Integration Step: \texttt{advance}}

\begin{Verbatim}[formatcom=\color{black}]
def advance(self, dt_s: float) -> dict[str, float]:
    """Integrate cable by dt_s seconds.
    
    Returns:
        dict mapping node names to activation times for nodes
        that crossed threshold during this step.
    """
    n_micro = max(1, round(dt_s / _DT_INT))
    dt_micro = dt_s / n_micro
    newly_activated = {}
    
    for _ in range(n_micro):
        self._t += dt_micro
        
        # Atrial and ventricular integration
        self._step_atrial(dt_micro)
        self._step_ventricular(dt_micro)
        
        # Detect threshold crossings (activation)
        for node_name, (i0, i1) in _ATR_SEGS.items():
            v_max = np.max(self._sa[i0:i1, 0])
            
            # If crossed threshold and not yet activated
            if v_max > V_THRESHOLD and self._t_act_a[i0] == math.inf:
                self._t_act_a[i0] = self._t  # Record activation time
                newly_activated[node_name] = self._t
        
        # Similar for ventricular nodes
        ...
    
    return newly_activated
\end{Verbatim}

\subsection{Numerical Integration Functions}

\begin{Verbatim}[formatcom=\color{black}]
def _step_atrial(self, dt: float) -> None:
    """One Euler step on atrial cable with stimulus."""
    V = self._sa[:, 0].copy()
    W = self._sa[:, 1].copy()
    
    # Stimulus current
    I = np.zeros(_NA)
    if self._t < self._stim_end_a:
        I[0] = _STIM_AMP  # Stimulus only at SA node
    
    # FHN step
    V_new, W_new = _fhn_cable_step(
        V, W, _D_ATR, I, dt,
        _TAU, _A, _B, _EPS
    )
    
    self._sa[:, 0] = V_new
    self._sa[:, 1] = W_new

@numba.njit(cache=True)  # JIT compile if Numba available
def _fhn_cable_step(V, W, D, I_ext, dt, tau, a, b, eps):
    """FHN dynamics on cable.
    
    Implements discretized equations (1)-(2) with finite-difference Laplacian.
    """
    N = len(V)
    
    # Laplacian (no-flux boundary conditions)
    lap = np.zeros(N)
    lap[1:-1] = V[:-2] - 2.0 * V[1:-1] + V[2:]
    lap[0]    = V[1] - V[0]
    lap[-1]   = V[-2] - V[-1]
    
    # FHN dynamics
    inv_tau = 1.0 / tau
    dv_dt = (D * lap + V - V**3 / 3.0 - W + I_ext) * inv_tau
    dw_dt = eps * (V + a - b * W) * inv_tau
    
    return V + dt * dv_dt, W + dt * dw_dt
\end{Verbatim}

Note: Numba JIT attempt occurs at import time; if incompatible NumPy version, falls back to pure NumPy (3–5$\times$ slower but functional).

\end{document}
